\section{Reduction II, part 1: the Jacobian at a root}\label{sec:jacobian}

\begin{lemma}\label{lem:jacobian}
  Fix $X\in\C^d$. Conditionally on $P(X)=0$, the rows of $J(X)$ are i.i.d.\ centered complex
  Gaussian vectors with covariance
  \[
    \Sigma(X) = S(X)\,H(X), \qquad
    H_{jk}(X) := \frac{\partial^2}{\partial x_j\,\partial\bar y_k}
    \log K(x,y)\Big|_{y=x=X};
  \]
  equivalently $J \overset{d}{=} G\,\Sigma(X)^{1/2}$ with $G$ as in
  Section~\ref{sec:background}.
\end{lemma}

\proof
Rows are independent because the $p_i$ are. For one row, apply Gaussian conditioning (Schur
complement) to the jointly Gaussian vector $(p_i(X),\partial_1p_i(X),\dots,\partial_dp_i(X))$:
with $K_j := \partial_{x_j}K|_X$ and $K_{j\bar k} := \partial_{x_j}\partial_{\bar y_k}K|_X$,
the conditional covariance is $K_{j\bar k} - K_j\overline{K_k}/S$, which equals
$S\,\partial_{x_j}\partial_{\bar y_k}\log K$ by expanding the logarithmic derivative.
\endproof

\begin{lemma}\label{lem:intensity}
  The expected number of roots per unit volume at $X$ is
  \[
    \rho(X) = \frac{d!}{\pi^d}\,\det H(X).
  \]
\end{lemma}

\proof
Apply \eqref{eq:counting} with $\varphi=1$:
$\rho(X) = \E[\,|\det J|^2 \mid P(X)=0\,]\,p_{P(X)}(0)$. Independence of the components gives
$p_{P(X)}(0) = (\pi S)^{-d}$. By Lemma~\ref{lem:jacobian} and $\E|\det G|^2 = d!$ (expand
$\det G\,\overline{\det G}$ over pairs of permutations; only equal pairs survive),
$\E[|\det J|^2 \mid P(X)=0] = d!\,S^d\det H$, and the powers of $S$ cancel.
\endproof

The factor $d!$ deserves verification; here are two independent checks.

\begin{proposition}\label{prop:masscheck}
  In the model of Definition~\ref{def:model}, $\int_{\C^d}\rho(X)\,dX = d!\,n^d = N$.
\end{proposition}

\proof
Since $\log K = \sum_l\log k_n(x_l\bar y_l)$, the matrix $H$ is diagonal with entries
$h(r_l)$ given in Lemma~\ref{lem:h}. In the radial variable $r=|z|^2$ one has $dA=\pi\,dr$,
so
\[
  \int_{\C^d}\det H\,\frac{dX}{\pi^d}
  = \prod_{l=1}^d\int_0^\infty h(r)\,dr
  = \prod_{l=1}^d\Bigl[\frac{r\,k_n'(r)}{k_n(r)}\Bigr]_0^\infty = n^d,
\]
because $h$ is the derivative of $rk_n'/k_n$, which tends to $\deg k_n = n$ at infinity. This
matches \eqref{eq:count}.
\endproof

\begin{proposition}\label{prop:kostlandensity}
  For the Kostlan kernel $K(x,y) = (1+\langle x,y\rangle)^n$, Lemma~\ref{lem:intensity} gives
  \[
    \rho(X) = \frac{d!\,n^d}{\pi^d(1+\|X\|^2)^{d+1}},
  \]
  which is $n^d$ times the projective root density derived independently in our earlier
  notes, with total mass equal to the B\'ezout number $n^d$.
\end{proposition}

\proof
Here $H = n\bigl[(1+s)^{-1}I - (1+s)^{-2}XX^*\bigr]$ with $s=\|X\|^2$; its eigenvalues are
$n/(1+s)$ with multiplicity $d-1$ and $n/(1+s)^2$, so $\det H = n^d(1+s)^{-(d+1)}$. The mass
integral $\int(1+\|X\|^2)^{-(d+1)}dX = \pi^d/d!$ is elementary.
\endproof
